\chapter{Coordinate Bethe Ansatz for the XXZ Chain}
\label{chap:coordinate-bethe-ansatz-xxz}

The preceding chapters introduced the XXZ chain as the Hamiltonian limit of the
six-vertex transfer matrix.  We now solve the same quantum chain directly in real
space.  This is the coordinate Bethe ansatz.  Its logic is deliberately close to
ordinary few-body quantum mechanics: write a plane-wave superposition in each
ordered coordinate sector, impose the Schrodinger equation at particle contact,
and finally impose periodic boundary conditions.

Throughout this chapter we use the integrability normalization obtained from the
six-vertex logarithmic derivative,
\begin{equation}
\label{eq:cba-xxz-hamiltonian}
    H_{\rm XXZ}^{(+)}
    =
    J\sum_{j=1}^{L}
    \left(
        \sigma_j^x\sigma_{j+1}^x
        +\sigma_j^y\sigma_{j+1}^y
        +\Delta\bigl(\sigma_j^z\sigma_{j+1}^z-1\bigr)
    \right),
\end{equation}
with periodic boundary conditions $\sigma_{L+1}^\alpha=\sigma_1^\alpha$.  The
fully polarized state
\begin{equation}
\label{eq:cba-ferromagnetic-reference}
    \ket{0}=\ket{\uparrow\uparrow\cdots\uparrow}
\end{equation}
is an eigenstate with zero energy.  It is not necessarily the ground state of the
antiferromagnetic chain, but it is the natural pseudovacuum for Bethe ansatz.
Down spins above this reference state will be called magnons.

The three goals of the chapter are:
\begin{enumerate}[label=(\roman*)]
    \item solve the one-magnon problem and identify the bare dispersion;
    \item derive the two-magnon scattering amplitude from the contact equation;
    \item assemble the $M$-magnon Bethe wave function and obtain the Bethe
    equations.
\end{enumerate}
The result is the first genuinely interacting exact solution in these notes: the
energy is additive in Bethe momenta, while the allowed momenta are determined by
many-body quantization conditions containing all two-body phase shifts.

\section{Conserved magnon number and coordinate basis}
\label{sec:cba-coordinate-basis}

Using
\begin{equation}
    \sigma_j^x\sigma_{j+1}^x+\sigma_j^y\sigma_{j+1}^y
    =
    2\left(\sigma_j^+\sigma_{j+1}^-+\sigma_j^-\sigma_{j+1}^+\right),
\end{equation}
one sees that the $XY$ part only exchanges neighboring up and down spins.  The
total magnetization
\begin{equation}
\label{eq:cba-total-magnetization}
    S^z_{\rm tot}=\frac12\sum_{j=1}^{L}\sigma_j^z
\end{equation}
therefore commutes with $H_{\rm XXZ}^{(+)}$.  Equivalently, the number of down
spins is conserved.

In the sector with $M$ down spins, a convenient basis is
\begin{equation}
\label{eq:cba-magnon-coordinate-basis}
    \ket{x_1,\ldots,x_M}
    =
    \sigma_{x_1}^-\cdots\sigma_{x_M}^-\ket{0},
    \qquad
    1\leq x_1<x_2<\cdots<x_M\leq L .
\end{equation}
A general state in this sector is
\begin{equation}
\label{eq:cba-general-M-magnon-state}
    \ket{\Psi_M}
    =
    \sum_{1\leq x_1<\cdots<x_M\leq L}
    \Psi(x_1,\ldots,x_M)\ket{x_1,\ldots,x_M} .
\end{equation}
The coordinate Bethe ansatz is an ansatz for the amplitudes
$\Psi(x_1,\ldots,x_M)$.

It is useful to separate two kinds of equations.  When no two magnons are
neighbors, the Schrodinger equation is just a discrete free-particle equation.
When $x_{a+1}=x_a+1$, the two magnons touch, and the nearest-neighbor Ising term
produces a contact condition.  This contact condition is the source of the
scattering phase.

\section{One-magnon sector}
\label{sec:cba-one-magnon-sector}

A one-magnon state has the form
\begin{equation}
\label{eq:cba-one-magnon-state}
    \ket{\Psi_1}
    =
    \sum_{x=1}^{L}\psi(x)\ket{x},
    \qquad
    \ket{x}=\sigma_x^-\ket{0} .
\end{equation}
Acting with \cref{eq:cba-xxz-hamiltonian}, one obtains
\begin{equation}
\label{eq:cba-one-magnon-difference-equation}
    E\psi(x)
    =
    2J\psi(x-1)+2J\psi(x+1)-4J\Delta\psi(x),
\end{equation}
where $x$ is understood modulo $L$.  The hopping amplitude is $2J$, and the
potential term $-4J\Delta$ comes from the two domain walls adjacent to the down
spin.

The plane wave
\begin{equation}
\label{eq:cba-one-magnon-plane-wave}
    \psi_k(x)=\ee^{\ii kx}
\end{equation}
is therefore an eigenfunction with dispersion
\begin{equation}
\label{eq:cba-one-magnon-dispersion}
    \varepsilon(k)
    =
    4J(\cos k-\Delta).
\end{equation}
Periodic boundary conditions give
\begin{equation}
\label{eq:cba-one-magnon-momentum-quantization}
    \ee^{\ii kL}=1,
    \qquad
    k=\frac{2\pi n}{L},
    \qquad
    n\in\ZZ .
\end{equation}

\begin{remark}[Energy convention]
The dispersion in \cref{eq:cba-one-magnon-dispersion} is written relative to the
fully polarized reference state, whose energy was fixed to zero by the constant
in \cref{eq:cba-xxz-hamiltonian}.  Reversing the overall sign of the Hamiltonian
or removing the constant shifts and rescales the formulas, but does not change
the Bethe wave functions or the scattering phase.
\end{remark}

\section{Two-magnon scattering}
\label{sec:cba-two-magnon-scattering}

The first nontrivial part of the Bethe ansatz appears for two magnons.  For
coordinates $x_1<x_2$, take
\begin{equation}
\label{eq:cba-two-magnon-wavefunction}
    \Psi(x_1,x_2)
    =
    A_{12}\,\ee^{\ii(k_1x_1+k_2x_2)}
    +
    A_{21}\,\ee^{\ii(k_2x_1+k_1x_2)} .
\end{equation}
The two terms describe the same two momenta assigned to the two ordered
coordinates in two different ways.  Away from contact, namely for
$x_2>x_1+1$, the wave function satisfies the free difference equation and has
energy
\begin{equation}
\label{eq:cba-two-magnon-energy}
    E(k_1,k_2)
    =
    4J(\cos k_1+\cos k_2-2\Delta).
\end{equation}
The energy is additive in the two bare one-magnon energies.

At contact, $x_2=x_1+1$, two of the four domain walls of separated magnons have
annihilated.  The Schrodinger equation becomes
\begin{equation}
\label{eq:cba-two-magnon-contact-equation}
    (E+4J\Delta)\Psi(x,x+1)
    =
    2J\Psi(x-1,x+1)+2J\Psi(x,x+2).
\end{equation}
Substituting \cref{eq:cba-two-magnon-wavefunction} into
\cref{eq:cba-two-magnon-contact-equation} fixes the ratio of amplitudes:
\begin{equation}
\label{eq:cba-two-magnon-scattering-amplitude}
    \frac{A_{21}}{A_{12}}
    =
    S(k_1,k_2)
    =
    -
    \frac{1+\ee^{\ii(k_1+k_2)}-2\Delta\ee^{\ii k_1}}
         {1+\ee^{\ii(k_1+k_2)}-2\Delta\ee^{\ii k_2}} .
\end{equation}
This is the two-magnon scattering amplitude of the XXZ chain in the convention
of \cref{eq:cba-xxz-hamiltonian}.

Several checks are immediate.  First,
\begin{equation}
\label{eq:cba-S-unitarity}
    S(k_1,k_2)S(k_2,k_1)=1,
\end{equation}
so exchanging two magnons twice gives back the original amplitude.  Second, at
$\Delta=0$,
\begin{equation}
    S(k_1,k_2)=-1,
\end{equation}
which is the free-fermion sign of the XX chain.  Thus the XX point is recovered
as the noninteracting limit of the coordinate Bethe ansatz.

\begin{principle}[Interaction as a boundary condition]
In the coordinate Bethe ansatz, the interaction does not destroy plane waves in
each ordered region.  Instead, it appears as a matching condition at contact,
which relates different plane-wave amplitudes by a two-body scattering phase.
\end{principle}

\section{The many-magnon coordinate Bethe ansatz}
\label{sec:cba-many-magnon-wavefunction}

For $M$ magnons, the ordered coordinate region is
\begin{equation}
    1\leq x_1<x_2<\cdots<x_M\leq L .
\end{equation}
The coordinate Bethe ansatz is
\begin{equation}
\label{eq:cba-M-magnon-wavefunction}
    \Psi(x_1,\ldots,x_M)
    =
    \sum_{P\in S_M}A(P)
    \exp\left(\ii\sum_{a=1}^{M}k_{P_a}x_a\right),
\end{equation}
where $S_M$ is the permutation group of $M$ objects.  The amplitude $A(P)$ is
attached to the assignment of momentum $k_{P_a}$ to the ordered coordinate
$x_a$.

Whenever all magnons are separated, the Schrodinger equation gives the additive
energy
\begin{equation}
\label{eq:cba-M-magnon-energy}
    E(k_1,
    \ldots,k_M)
    =
    \sum_{a=1}^{M}\varepsilon(k_a)
    =
    4J\sum_{a=1}^{M}(\cos k_a-\Delta).
\end{equation}
Whenever two adjacent coordinates touch, the same two-body contact equation as in
the two-magnon problem applies.  Therefore adjacent permutations obey
\begin{equation}
\label{eq:cba-adjacent-permutation-relation}
    A(Ps_a)
    =
    S(k_{P_a},k_{P_{a+1}})A(P),
    \qquad
    a=1,\ldots,M-1,
\end{equation}
where $s_a$ exchanges the entries in positions $a$ and $a+1$.

Because the scattering amplitude is a scalar phase, consistency of different
ways of reducing a permutation follows from ordinary multiplication.  In models
with internal degrees of freedom, the analogous consistency condition becomes the
Yang--Baxter equation.  This is why the coordinate Bethe ansatz here is the
real-space shadow of the six-vertex Yang--Baxter structure.

\section{Periodic boundary conditions and Bethe equations}
\label{sec:cba-bethe-equations-momentum-form}

The remaining step is to impose periodic boundary conditions.  Move magnon $j$
once around the ring.  It accumulates a free propagation phase $\ee^{\ii k_jL}$
and scatters once with every other magnon.  With the scattering convention of
\cref{eq:cba-two-magnon-scattering-amplitude}, the result is
\begin{equation}
\label{eq:cba-momentum-bethe-equations}
    \ee^{\ii k_jL}
    =
    \prod_{\substack{\ell=1\\ \ell\neq j}}^{M}
    S(k_j,k_\ell),
    \qquad
    j=1,\ldots,M .
\end{equation}
Together with the energy formula \cref{eq:cba-M-magnon-energy}, these equations
solve the finite periodic XXZ chain in each fixed-magnon sector, up to the usual
questions of completeness and string organization.

Writing
\begin{equation}
    S(k_j,k_\ell)=\ee^{\ii\theta(k_j,k_\ell)},
\end{equation}
one obtains the logarithmic form
\begin{equation}
\label{eq:cba-logarithmic-bethe-equations}
    Lk_j
    =
    2\pi I_j
    +
    \sum_{\substack{\ell=1\\ \ell\neq j}}^{M}
    \theta(k_j,k_\ell),
    \qquad
    I_j\in\ZZ\ \text{or}\ \ZZ+\frac12 .
\end{equation}
The integer versus half-integer choice depends on $M$, $L$, the branch choice of
$\theta$, and the boundary-condition convention.  The exponential form
\cref{eq:cba-momentum-bethe-equations} is the unambiguous statement.

\begin{remark}[What is solved by the Bethe equations]
The Bethe equations do not say that the interacting chain has free momenta.  They
say that eigenstates are still labelled by rapidities or momenta, but those
labels are quantized collectively.  Each $k_j$ is shifted by the phase shifts
against all other magnons.
\end{remark}

\section{Rapidity parametrizations}
\label{sec:cba-rapidity-parametrizations}

The momentum-space scattering amplitude is explicit but not the most useful form
for analysis.  The standard Bethe equations become simpler after introducing a
rapidity variable.  The appropriate parametrization depends on the anisotropy
regime.

\subsection*{Gapless regime: $-1<\Delta<1$}

Let
\begin{equation}
\label{eq:cba-gapless-parametrization}
    \Delta=\cos\gamma,
    \qquad
    0<\gamma<\pi .
\end{equation}
Define the rapidity $\lambda$ by
\begin{equation}
\label{eq:cba-gapless-momentum-rapidity-map}
    \ee^{\ii k(\lambda)}
    =
    \frac{\sinh(\lambda+\ii\gamma/2)}
         {\sinh(\lambda-\ii\gamma/2)} .
\end{equation}
Then \cref{eq:cba-momentum-bethe-equations} becomes
\begin{equation}
\label{eq:cba-gapless-xxz-bethe-equations}
    \left[
    \frac{\sinh(\lambda_j+\ii\gamma/2)}
         {\sinh(\lambda_j-\ii\gamma/2)}
    \right]^L
    =
    \prod_{\substack{\ell=1\\ \ell\neq j}}^{M}
    \frac{\sinh(\lambda_j-\lambda_\ell+\ii\gamma)}
         {\sinh(\lambda_j-\lambda_\ell-\ii\gamma)} .
\end{equation}
The energy becomes
\begin{equation}
\label{eq:cba-gapless-energy-rapidity}
    E
    =
    -4J\sum_{j=1}^{M}
    \frac{\sin^2\gamma}{\cosh(2\lambda_j)-\cos\gamma} .
\end{equation}
The total lattice momentum is
\begin{equation}
\label{eq:cba-total-momentum}
    P
    =
    \sum_{j=1}^{M}k(\lambda_j)
    \quad (\operatorname{mod} 2\pi).
\end{equation}

\subsection*{Massive antiferromagnetic regime: $\Delta>1$}

Let
\begin{equation}
\label{eq:cba-massive-parametrization}
    \Delta=\cosh\eta,
    \qquad
    \eta>0 .
\end{equation}
A convenient rapidity map is
\begin{equation}
\label{eq:cba-massive-momentum-rapidity-map}
    \ee^{\ii k(\lambda)}
    =
    \frac{\sin(\lambda+\ii\eta/2)}
         {\sin(\lambda-\ii\eta/2)} .
\end{equation}
The Bethe equations are
\begin{equation}
\label{eq:cba-massive-xxz-bethe-equations}
    \left[
    \frac{\sin(\lambda_j+\ii\eta/2)}
         {\sin(\lambda_j-\ii\eta/2)}
    \right]^L
    =
    \prod_{\substack{\ell=1\\ \ell\neq j}}^{M}
    \frac{\sin(\lambda_j-\lambda_\ell+\ii\eta)}
         {\sin(\lambda_j-\lambda_\ell-\ii\eta)} .
\end{equation}
The energy is
\begin{equation}
\label{eq:cba-massive-energy-rapidity}
    E
    =
    -4J\sum_{j=1}^{M}
    \frac{\sinh^2\eta}{\cosh\eta-\cos(2\lambda_j)} .
\end{equation}
In this regime complex rapidities organize into string patterns in the
thermodynamic limit.  Strings describe bound states of magnons and are essential
for thermodynamics, but the finite-size Bethe equations above are the starting
point.

\subsection*{Isotropic XXX point}

At $\Delta=1$, one may use the rational rapidity variable
\begin{equation}
\label{eq:cba-xxx-momentum-rapidity-map}
    \ee^{\ii k(\lambda)}
    =
    \frac{\lambda+\ii/2}{\lambda-\ii/2} .
\end{equation}
The Bethe equations reduce to
\begin{equation}
\label{eq:cba-xxx-bethe-equations}
    \left(
    \frac{\lambda_j+\ii/2}{\lambda_j-\ii/2}
    \right)^L
    =
    \prod_{\substack{\ell=1\\ \ell\neq j}}^{M}
    \frac{\lambda_j-\lambda_\ell+\ii}
         {\lambda_j-\lambda_\ell-\ii} ,
\end{equation}
and the energy is
\begin{equation}
\label{eq:cba-xxx-energy}
    E
    =
    -2J\sum_{j=1}^{M}
    \frac{1}{\lambda_j^2+1/4} .
\end{equation}
This is the rational limit of the trigonometric XXZ equations.

\section{Special limits and physical interpretation}
\label{sec:cba-special-limits-interpretation}

\begin{example}[XX point]
At $\Delta=0$, the scattering amplitude becomes $S=-1$.  The Bethe equations are
therefore
\begin{equation}
    \ee^{\ii k_jL}=(-1)^{M-1} .
\end{equation}
This is precisely the parity-dependent boundary condition of Jordan--Wigner
fermions in the $M$-particle sector.  Thus the coordinate Bethe ansatz reduces to
free fermions at the XX point.
\end{example}

\begin{example}[Two-magnon bound states]
For $\Delta>1$, the scattering amplitude has poles in the complex momentum or
rapidity plane.  Normalizable two-magnon wave functions can then have complex
conjugate momenta, producing an exponentially localized relative coordinate.
These states are the simplest examples of Bethe strings.  In the many-body
problem, strings encode bound-state quasiparticles.
\end{example}

\begin{remark}[Pseudovacuum versus physical ground state]
The coordinate Bethe ansatz above builds states by flipping spins above
$\ket{0}=\ket{\uparrow\cdots\uparrow}$.  For the antiferromagnetic convention
with $J>0$ and $\Delta>0$, this reference state is not the physical ground state
in the zero-magnetization sector.  The true ground state is obtained by solving
the Bethe equations with an extensive number of rapidities.  The pseudovacuum is
used because it makes the scattering construction algebraically simple.
\end{remark}

\section{Relation to the algebraic Bethe ansatz}
\label{sec:cba-relation-to-aba}

The coordinate Bethe ansatz diagonalizes the Hamiltonian by solving the real-space
Schrodinger equation.  The algebraic Bethe ansatz diagonalizes the whole transfer
matrix family using the $R$-matrix algebra.  The two methods produce the same
Bethe equations, but they organize the calculation differently:
\begin{center}
\begin{tabular}{@{}L{0.30\linewidth}L{0.30\linewidth}L{0.30\linewidth}@{}}
\toprule
Method & Basic object & Origin of Bethe equations \\
\midrule
Coordinate Bethe ansatz & Wave function
$\Psi(x_1,\ldots,x_M)$ & Periodic boundary conditions plus two-body scattering \\
Algebraic Bethe ansatz & Bethe vector
$B(\lambda_1)\cdots B(\lambda_M)\ket{0}$ & Cancellation of unwanted terms in the RTT algebra \\
\bottomrule
\end{tabular}
\end{center}
The coordinate method is more physically transparent: it shows directly that the
XXZ chain is solvable because many-body scattering factorizes into two-body
phases.  The algebraic method is more structural: it explains why this
factorization is guaranteed by the Yang--Baxter equation and gives direct access
to transfer-matrix eigenvalues and higher conserved charges.

\section{Summary}
\label{sec:cba-summary}

The coordinate Bethe ansatz for the XXZ chain can be summarized as
\begin{equation}
\label{eq:cba-summary-chain}
\boxed{
\begin{gathered}
    H_{\rm XXZ}^{(+)}
    \quad\Longrightarrow\quad
    M\text{-magnon coordinate sectors}
    \\
    \Longrightarrow\quad
    \Psi(x_1,\ldots,x_M)
    =
    \sum_{P\in S_M}A(P)
    \exp\left(\ii\sum_a k_{P_a}x_a\right)
    \\
    \Longrightarrow\quad
    A(Ps_a)=S(k_{P_a},k_{P_{a+1}})A(P)
    \\
    \Longrightarrow\quad
    \ee^{\ii k_jL}=\prod_{\ell\neq j}S(k_j,k_\ell),
    \qquad
    E=4J\sum_j(\cos k_j-\Delta).
\end{gathered}}
\end{equation}
The one-magnon sector gives the bare dispersion.  The two-magnon contact equation
gives the scattering phase.  The many-magnon solution works because all
collisions factorize into the same two-body scattering amplitude.  This is the
coordinate-space form of the Yang--Baxter integrability already encountered in
the six-vertex model.
